\chapter{Conclusion}
\label{sec:conclusion}

This thesis investigated five questions concerning model comparison, threshold selection, imbalance interventions, BAF Base-to-Variant transfer, and model interpretation in fraud detection under extreme class imbalance. The experimental programme covers two benchmark datasets, six model families, seven imbalance handling strategies for the five supervised models, four threshold rules, SHAP analysis for compatible model families, and an FT-Transformer attention diagnostic. A shared leakage-free outer holdout protects the final predictive evaluation, while the model-specific inner-validation procedures remain explicit.

% ──────────────────────────────────────────────────────────────────────────
\section{Summary of Contributions}
\label{sec:contributions_summary}
% ──────────────────────────────────────────────────────────────────────────

The six contributions follow the sequence established by the research design; Contributions~5 and~6 jointly address RQ5.

The foundation is a \textbf{leakage-free unified comparison} (Contribution~1) that provides a side-by-side evaluation of classical supervised models and FT-Transformer on the same two datasets, with common metrics and an outer holdout that is excluded from preprocessing, model selection, resampling, and threshold calibration. The classical models use five-fold inner validation and FT-Transformer uses a single internal holdout, so the comparison is unified at the outer-evaluation level rather than identical at every inner-training step.

Built on that foundation, the \textbf{threshold selection study} (Contribution~2) shows that the decision rule materially changes threshold-dependent metrics. On BAF Base, fixing the threshold at 0.5 yields $F_2<0.05$ for every supervised model, whereas validation-selected max-$F_2$ thresholds produce values in the 0.29--0.32 range. For LGBM, CatBoost, and FT-Transformer, this corresponds to approximately 7.5-, 8.5-, and 9.1-fold increases, respectively. Within the evaluated workflow, this result supports explicit validation-based threshold selection whenever operational decisions require a binary cut-off.

The \textbf{full factorial analysis} (Contribution~3) across seven imbalance handling strategies, five supervised models, and two datasets shows that strategy effectiveness is neither universal nor model-agnostic. On BAF, SMOTE reduces FT-Transformer PR-AUC by 41\% (0.180 to 0.106) and CatBoost by 14\% (0.180 to 0.154). However, the FT-Transformer sampler interpolates ordinal category codes before converting the synthetic values to integer embedding indices. The result therefore shows sensitivity under the implemented resampling representation, not a structural weakness attributable solely to neural architecture. Class Weights remains close to the FT-Transformer baseline on BAF but reduces its ULB PR-AUC, so no intervention is supported as a universal default.

The \textbf{cross-domain generalisation study} (Contribution~4) shows that Base performance alone does not describe the Variant profile. Random Forest records the highest PR-AUC and $F_2$ point estimates among the evaluated models on all five Variant test sets. Relative to its own Base result, its PR-AUC is higher on Variants~I--IV and effectively unchanged on Variant~V, although its $F_2$ decreases on Variants~III and~V. LGBM, CatBoost, and FT-Transformer show their largest PR-AUC reductions on Variants~III and~V, which alter group-wise and temporal separability. These descriptive comparisons use one Base-trained model and the unchanged Base threshold; no in-domain Variant models were trained, and PR-AUC differences may also reflect prevalence changes.

The \textbf{SHAP cross-model and cross-variant analysis} (Contribution~5) records two bounded results. LGBM and CatBoost have the same top-10 feature set (Jaccard\,=\,1.00), and each shares eight of ten members with Logistic Regression (Jaccard\,=\,0.67). For the fixed Base-trained LGBM model, the same top-10 set is recovered on all six BAF test distributions. These results concern set membership, not identical ordering, attribution magnitude, local-explanation validity, or causal importance. FT-Transformer GradientExplainer results are excluded from this contribution because integer conversion of categorical indices prevents reliable gradient attribution to those inputs.

Finally, the \textbf{FT-Transformer attention diagnostic} (Contribution~6) finds diffuse final-layer, head-averaged \texttt{[CLS]} attention (normalised entropy\,=\,0.976) across the 31 tokens. This result describes one aggregated attention view; it does not establish an information ceiling, show that all signal has been used, or rule out head-, layer-, or architecture-specific limitations.

% ──────────────────────────────────────────────────────────────────────────
\section{Answers to Research Questions}
\label{sec:rq_answers}
% ──────────────────────────────────────────────────────────────────────────

\paragraph{RQ1 — Comparative performance under a leakage-free protocol.}
Under the shared leakage-free outer evaluation, CatBoost has the highest ULB PR-AUC (0.885), followed by LGBM (0.874), RF (0.867), FT-Transformer (0.856), Logistic Regression (0.742), and OCSVM (0.335). On BAF Base, FT-Transformer, CatBoost, and LGBM are closely grouped at 0.180, 0.180, and 0.177 by rounded PR-AUC. The observed gap is therefore dataset-dependent, but the experiments do not isolate dataset size, calibration, or inductive bias as its cause. FT-Transformer does not exceed the highest GBDT point estimate on either benchmark and requires substantially greater measured compute than LGBM (Section~\ref{sec:baseline_results}, p.~\pageref{sec:baseline_results}).

\paragraph{RQ2 — Impact of threshold selection strategy.}
Threshold selection has a dataset- and model-dependent impact on operational metrics. ULB changes are moderate for most models but reach a 20\% relative $F_2$ increase for Logistic Regression. On BAF, fixed 0.5 gives $F_2<0.05$ for every supervised model, while max-$F_2$ produces 0.29--0.32. Rankings change under different rules on both datasets, so any claim about the best model under a threshold-dependent metric must specify the operating rule. All non-fixed thresholds are selected inside DEV, and the precision constraint is a validation target rather than a TEST guarantee (Section~\ref{sec:threshold_sensitivity}, p.~\pageref{sec:threshold_sensitivity}).

\paragraph{RQ3 — Factorial interaction of imbalance strategies and model type.}
No single imbalance strategy improves every model and dataset. The baseline is the highest-PR-AUC condition for four of five supervised models on ULB, while ROS gives a small Random Forest improvement. On BAF, LGBM and CatBoost gain less than 0.01 PR-AUC from any intervention. FT-Transformer's 41\% SMOTE reduction is larger than CatBoost's 14\% reduction, but the categorical resampling representation is a confound. Class Weights is close to the FT-Transformer baseline on BAF and lower on ULB; strategy selection should therefore be validated for the target model and data rather than chosen as a universal default (Sections~\ref{sec:factorial_results} and~\ref{sec:transformer_imbalance}, pp.~\pageref{sec:factorial_results} and~\pageref{sec:transformer_imbalance}).

\paragraph{RQ4 — Cross-domain generalisation under distribution shift.}
Base performance ranking does not determine every BAF Variant point estimate. Random Forest records the highest PR-AUC and $F_2$ point estimates among the evaluated models on each of the five Variant test sets. Variant~II, which changes group-wise prevalence disparity while equalising group sizes, is close to or above Base for several models. Variants~III and~V, which alter separability, produce the largest PR-AUC reductions for LGBM, CatBoost, and FT-Transformer. Relative to its own Base result, Random Forest is higher on Variants~I--IV by PR-AUC and essentially unchanged on Variant~V, but its $F_2$ decreases on Variants~III and~V. The study does not isolate the mechanism behind this profile; robustness should be evaluated across the relevant shifts and metrics rather than inferred from Base performance alone (Section~\ref{sec:cross_domain_results}, p.~\pageref{sec:cross_domain_results}).

\paragraph{RQ5 — Feature attribution consistency and transformer attention diagnostics.}
LGBM and CatBoost have the same top-10 SHAP feature set, and each overlaps Logistic Regression on eight members. The fixed Base-trained LGBM model also retains the same top-10 set across all BAF test distributions. These are set-membership results and do not establish identical rankings or explanation validity under shift. FT-Transformer SHAP is excluded from cross-model conclusions because of its categorical-gradient limitation. Its final-layer, head-averaged attention is diffuse (normalised entropy\,=\,0.976), which characterises the inspected attention view but does not prove an information limit or explain predictive parity (Sections~\ref{sec:shap_results} and~\ref{sec:attention_diagnostics}, pp.~\pageref{sec:shap_results} and~\pageref{sec:attention_diagnostics}).

% ──────────────────────────────────────────────────────────────────────────
\section{Limitations}
\label{sec:limitations}
% ──────────────────────────────────────────────────────────────────────────

Several limitations constrain the scope and generalisability of the findings beyond the protocol-level threats to validity discussed in Section~\ref{sec:discussion} (p.~\pageref{sec:discussion}).

\textbf{Single deep learning architecture.} FT-Transformer is the sole representative of the tabular deep learning family. The conclusions about attention behaviour, imbalance sensitivity, and cross-domain performance apply directly to this architecture and may not generalise to SAINT \cite{somepalli2021saint}, TabTransformer \cite{huang2020tabtransformer}, TabNet, or hybrid architectures. Whether similarly diffuse attention occurs in other architectures remains an open question.

\textbf{Model-specific validation.} The classical models use five-fold inner validation, whereas FT-Transformer uses one internal DEV holdout. Both paths protect the outer TEST set, but they do not provide equivalent estimates of selection and threshold variability. The fixed split and random seed also leave training variability unmeasured.

\textbf{Resampling representation.} The implemented SMOTE path is not categorical-aware. It interpolates one-hot columns for the classical models and ordinal category codes for FT-Transformer, whose synthetic categorical values are then converted to integers. The observed FT-Transformer reduction under SMOTE therefore cannot be attributed solely to architecture.

\textbf{Interpretability scope.} FT-Transformer GradientExplainer cannot reliably attribute categorical embedding indices because of the integer conversion in the explainer wrapper, so those results are excluded from cross-model conclusions. The attention analysis covers the final layer and averages across heads; it is descriptive and cannot establish feature causality or an information ceiling.

\textbf{Dataset scope.} The ULB dataset is privacy-constrained — PCA-transformed features preclude domain interpretation and the two-day collection window limits temporal generalisability. The BAF suite is synthetically generated, which ensures experimental control but means fraud patterns may not fully reflect the complexity of real production fraud streams, where label noise, adversarial adaptation, and operational feedback loops are present.

\textbf{Absence of concept drift modelling.} All experiments use a static train/test split. In operational fraud detection, the data distribution may shift as fraud patterns and system use evolve. The experimental protocol does not capture this dynamic, and performance under temporal drift is not equivalent to performance on the controlled BAF Variants studied here.

\textbf{Hyperparameter optimisation budget.} The same 50-trial count was applied to each supervised model, but OCSVM was not tuned and the validation cost per trial differs between five-fold classical searches and the FT-Transformer holdout. The FT-Transformer's larger search space may remain under-explored relative to the gradient-boosted models.

\textbf{Cross-dataset incompatibility.} The ULB and BAF feature spaces are incompatible (PCA-anonymised vs.\ semantic), preventing direct transfer between the datasets. The cross-domain study is therefore restricted to BAF internal transfer, a narrower robustness evaluation than matched feature spaces across independent datasets would provide.

% ──────────────────────────────────────────────────────────────────────────
\section{Future Work}
\label{sec:future_work}
% ──────────────────────────────────────────────────────────────────────────

Five directions emerge directly from the findings and limitations of this thesis.

\textbf{Extending attention diagnostics to other tabular transformer architectures.} The diffuse final-layer average observed for FT-Transformer motivates analysis of individual heads, earlier layers, and alternative tabular transformers. Applying a comparable diagnostic to SAINT and TabTransformer could assess whether attention concentration differs by architecture and whether any difference is associated with predictive behaviour.

\textbf{Temporal robustness and concept drift.} The BAF dataset preserves the temporal ordering of bank account applications, enabling time-ordered train/test splits that approximate deployment drift. A direct extension would compare periodic threshold recalibration with full model retraining as the distribution evolves, since the two approaches have different operational costs.

\textbf{Ensemble strategies based on predictive diversity.} LGBM and FT-Transformer reach similar BAF aggregate performance, but the current SHAP implementation cannot establish complementary FT-Transformer feature use. Future work could measure error and score correlation directly before testing whether a stacking ensemble offers predictive benefit.

\textbf{Native cost-sensitive training for the FT-Transformer.} Class Weights remains close to baseline on BAF but reduces performance on ULB, so its effect is dataset-dependent. Future work could compare focal loss or direct $F_{\beta}$ optimisation with categorical-aware resampling under a protocol designed specifically to separate loss-level and representation-level effects.

\textbf{Richer feature representations.} Augmenting BAF with derived behavioural features, velocity-based aggregates, graph connectivity, or device-fingerprinting signals would test whether the relative performance of transformer and tree models changes when additional cross-feature structure is available. The present attention diagnostic does not determine in advance whether such augmentation would help.

% ──────────────────────────────────────────────────────────────────────────
\section{Closing Remarks}
% ──────────────────────────────────────────────────────────────────────────

The central conclusion is methodological as well as predictive: model choice cannot be separated from validation design, threshold selection, imbalance handling, input representation, and the distribution on which the model is evaluated. Within the scope of the two benchmarks, gradient-boosted trees provide strong performance at substantially lower measured computational cost than FT-Transformer, while the transformer remains competitive on BAF Base. These findings support evidence-bounded model selection rather than a universal preference for either model family.
